\documentclass[t,aspectratio=169,10pt]{beamer}

\usepackage[utf8]{inputenc}

\usepackage{multicol}
\usepackage{multirow}
\usepackage{textcomp} % to use \textmu
\usepackage[absolute,overlay]{textpos} % to place floating text boxes with \begin{textblock*}{width}(x,y)
\usepackage{tcolorbox}

\usepackage{tikz}
\usetikzlibrary{decorations.pathreplacing}

\usepackage{graphicx}
\graphicspath{{./figuras/}}

\setbeamertemplate{navigation symbols}{}

\title{Clase 14}
\subtitle{Modelo de CD del transistor MOSFET}
\author[Dr.-Ing. Juan José Montero-Rodríguez]{Dr.-Ing. Juan José Montero-Rodríguez}
\subject{Elementos Activos}
\institute{Escuela de Ingeniería Electrónica}
\date{Semestre I-2024}
\titlegraphic{\includegraphics[height=8mm]{logotec.pdf}}


\begin{document}

\begin{frame}
\titlepage
\end{frame}


\begin{frame}{Modelos del MOSFET}
\centering
\includegraphics[width=11cm]{modelosMOSFET}
\end{frame}


\begin{frame}{Modelo Analógico del MOSFET}
Modelo analógico según región de operación

\centering\vspace{3mm}
\includegraphics[width=11cm]{analogico1}
\end{frame}



\section{Curvas}
\begin{frame}{Curvas Características del MOSFET: Curva de Transferencia}

La curva de transferencia se mide al dejar $V_{DS}$ constante y modificar $V_{GS}$.

\begin{figure}
    \centering
    \includegraphics[width=0.8\textwidth]{figuras/mosfet_transfer.png}
\end{figure}



\end{frame}



\begin{frame}{Curvas Características del MOSFET: Curva de Salida}

La curva de salida se mide al dejar $V_{GS}$ constante y modificar $V_{DS}$.

\begin{figure}
    \centering
    \includegraphics[width=0.8\textwidth]{figuras/mosfet_output.png}
\end{figure}\textbf{}

\end{frame}


\section{Regiones}
\begin{frame}{Región lineal}

\begin{columns}

	\begin{column}{0.5\textwidth}
 
		\centering
		\textbf{Región Lineal}
		%
		\[ V_{GS} > V_{TH}, V_{DS} < V_{GS} - V_{TH} \]
		%
		\includegraphics[width=5cm]{lineal}
        
        El canal está completo

        \vspace{3mm}
        Funcionamiento como resistencia controlada por tensión
  
	\end{column}
 
	\begin{column}{0.5\textwidth}
 
        \begin{figure}
            \centering
            \includegraphics[width=\textwidth]{figuras/region_lineal.png}
        \end{figure}

        \[ I_D = \mu C_{ox}' \dfrac{W}{L} \left( (V_{GS} - V_{TH}) V_{DS} - \dfrac{V_{DS}^2}{2} \right) \]
  
	\end{column}
 
\end{columns}

\end{frame}


\begin{frame}{Región lineal: derivación de ecuaciones}

\begin{columns}

\begin{column}{0.5\textwidth}

La carga eléctrica $[C/m]$:
\[ Q = C \cdot V \]

La corriente eléctrica $[A]$:
\[ I = Q \cdot v \]
%
\[ I = (C V) \cdot (\mu E) \]

Para el MOSFET se tiene:
\[ C = W C_{ox}' \]

Combinando ecuaciones:
\[ I =  W C_{ox}' V \mu \dfrac{dV}{dx} \]

\end{column}

\begin{column}{0.5\textwidth}

La tensión es función de la posición:
\[ I = C \left( V_{GS} - V(x) - V_{TH} \right) \mu \dfrac{dV}{dx} \]
%
\[ I \cdot dx = C (V_{GS} - V(x) - V_{TH}) \mu dV \]
%
\[ \int_0^L I \cdot dx = \int_0^{V_{DS}} C (V_{GS} - V(x) - V_{TH}) \mu dV \]
%
\[ I \cdot L = \mu C_{ox}' W \left( (V_{GS} - V_{TH})V_{DS} - \dfrac{V_{DS}^2}{2} \right) \]
%
\[ \boxed{I = \mu C_{ox}' \dfrac{W}{L} \left( (V_{GS} - V_{TH})V_{DS} - \dfrac{V_{DS}^2}{2} \right)} \]
%
\centering Corriente de drenador en región lineal

\end{column}
    
\end{columns}


    
\end{frame}



\begin{frame}{Región lineal: resistencia de canal}
MOSFET se comporta como resistencia R\textsubscript{ch} (resistencia de canal)
%
\[ R_{ch} = \left( \dfrac{\partial I_D}{\partial V_{DS}} \right)^{-1} \]
%
\[ \dfrac{\partial I_D }{\partial V_{DS}} = \mu C_{ox}' \dfrac{W}{L} (V_{GS} - V_{TH} - V_{DS})\]
%
Por lo tanto, la resistencia de canal se calcula como
%
\[ R_{ch} = \dfrac{1}{\mu C_{ox}' \dfrac{W}{L} (V_{GS} - V_{TH} - V_{DS})} \]
%
lo cual para $V_{DS} << (V_{GS} - V_{TH})$ puede aproximarse a
%
\[ \boxed{R_{ch} \approx \dfrac{1}{\mu C_{ox}' \dfrac{W}{L} (V_{GS} - V_{TH})}} \]
\end{frame}



\begin{frame}{Región de saturación}

\begin{columns}

	\begin{column}{0.5\textwidth}

		\centering
		\textbf{Región de Saturación}
		%
		\[ V_{GS} > V_{TH}, V_{DS} \geq V_{GS} - V_{TH} \]
		%
		\includegraphics[width=5cm]{satura}
		
		Estrangulamiento del canal
		
		\vspace{3mm}
		Funcionamiento como fuente de corriente controlada por tensión
 
	\end{column}
 
	\begin{column}{0.5\textwidth}
 
        \begin{figure}
            \centering
            \includegraphics[width=\textwidth]{figuras/region_saturacion.png}
        \end{figure}

        \[ I_D = \dfrac{1}{2} \mu {}C_{ox}' \dfrac{W}{L} (V_{GS}-V_{TH})^2 \]
  
	\end{column}
 
\end{columns}

\end{frame}


\begin{frame}{Región de saturación: derivación de ecuaciones}

La corriente de drenador en la región lineal:
\[ I = \mu C_{ox}' \dfrac{W}{L} \left( (V_{GS} - V_{TH})V_{DS} - \dfrac{V_{DS}^2}{2} \right) \]

El punto máximo de esta ecuación se encuentra igualando la primera derivada a cero:
\[ \dfrac{d}{dV_{DS}} \left\{ \mu C_{ox}' \dfrac{W}{L} \left( (V_{GS} - V_{TH})V_{DS} - \dfrac{V_{DS}^2}{2} \right) \right\} = 0 \]
%
\[ \mu C_{ox}' \dfrac{W}{L} \left( (V_{GS} - V_{TH}) - V_{DS} \right) = 0 \]
%
\[ V_{GS} - V_{TH} - V_{DS} = 0 \]
%
\[ V_{GS} - V_{TH} = V_{DS} \]

Por lo tanto, sustituyendo $V_{DS}$ en la ecuación para encontrar el máximo, se obtiene:
\[ \boxed{I = \dfrac{1}{2} \mu C_{ox}' \dfrac{W}{L} (V_{GS} - V_{TH})^2} \]
  
\end{frame}


\begin{frame}{Región de saturación: funcionamiento como fuente de corriente}
MOSFET se comporta como fuente de corriente en paralelo con resistencia

\centering
\includegraphics[width=11cm]{analogico2}

\includegraphics[width=11cm]{analogico3}
\end{frame}



\begin{frame}{Ecuaciones características del MOSFET}

Corriente de drenador en la región lineal: $V_{GS} > V_{TH}$, $V_{DS} < V_{GS} - V_{TH}$

\begin{tcolorbox}
\[ I_D = \mu C_{ox}' \dfrac{W}{L} \left( (V_{GS} - V_{TH}) V_{DS} - \dfrac{V_{DS}^2}{2} \right) \]
\end{tcolorbox}

Corriente de drenador en la región de saturación: $V_{GS} > V_{TH}$, $ V_{DS} > V_{GS} - V_{TH} $

\begin{tcolorbox}
\[ I_D = \dfrac{1}{2} \mu {}C_{ox}' \dfrac{W}{L} (V_{GS}-V_{TH})^2 \]
\end{tcolorbox}

\begin{columns}
	\begin{column}{0.25\textwidth}
		\centering Transconductancia del proceso:
		%
		\[ K' = \mu\cdot{}C_{ox}'\ [A/V^2] \] 
	\end{column}
	\begin{column}{0.25\textwidth}
		\centering Parámetro de transconductancia:
		%
		\[ K = \mu\cdot{}C_{ox}' \cdot \dfrac{W}{L} \]
	\end{column}
	\begin{column}{0.25\textwidth}
        \centering Capacitancia de óxido:
        %
        \[ C_{ox}' = \dfrac{\epsilon_{ox}}{t_{ox}}\ [F/m^2] \]
	\end{column}
	\begin{column}{0.25\textwidth}		
		$\mu$: movilidad
		
		$\epsilon_{ox}$: 3.9$\epsilon_0$
		
		$t_{ox}$: espesor de óxido
	\end{column}

\end{columns}
\end{frame}


\section{Ejemplos}
\begin{frame}{Ejemplo 1: MOSFET en saturación}

Determine el punto de operación de M1. Asuma $\mu_n C_{ox}' = 100\ \mu{}A/V^2$ y $V_{TH} = 0.4\ V$. Si la tensión de la compuerta aumenta 10 mV, ¿cuál es el cambio en la corriente de drenador?

\begin{columns}

\begin{column}{0.5\textwidth}

\begin{figure}[H]
    \centering
    \includegraphics[width=0.7\textwidth]{figuras/mosfet_ejemplo_1.png}
\end{figure}

\end{column}

\begin{column}{0.5\textwidth}

\end{column}

\end{columns}

\end{frame}


\begin{frame}{Solución 1: MOSFET en saturación}

Determine el punto de operación de M1. Asuma $\mu_n C_{ox}' = 100\ \mu{}A/V^2$ y $V_{TH} = 0.4\ V$. Si la tensión de la compuerta aumenta 10 mV, ¿cuál es el cambio en la corriente de drenador?

\begin{columns}

\begin{column}{0.5\textwidth}

\begin{figure}[H]
    \centering
    \includegraphics[width=0.7\textwidth]{figuras/mosfet_ejemplo_1.png}
\end{figure}

\end{column}

\begin{column}{0.5\textwidth}

\end{column}

\end{columns}

\end{frame}


\begin{frame}{Ejemplo 2: MOSFET en saturación}

Calcule la máxima tensión de compuerta permitida si se quiere que M1 permanezca en saturación. Asuma $\mu_n C_{ox}' = 100\ \mu{}A/V^2$ y $V_{TH} = 0.4\ V$.

\begin{columns}

\begin{column}{0.5\textwidth}

\begin{figure}[H]
    \centering
    \includegraphics[width=0.7\textwidth]{figuras/mosfet_ejemplo_2.png}
\end{figure}

\end{column}

\begin{column}{0.5\textwidth}

\end{column}

\end{columns}
    
\end{frame}


\begin{frame}{Solución 2: MOSFET en saturación}

Calcule la máxima tensión de compuerta permitida si se quiere que M1 permanezca en saturación. Asuma $\mu_n C_{ox}' = 100\ \mu{}A/V^2$ y $V_{TH} = 0.4\ V$.

\begin{columns}

\begin{column}{0.5\textwidth}

\begin{figure}[H]
    \centering
    \includegraphics[width=0.7\textwidth]{figuras/mosfet_ejemplo_2.png}
\end{figure}

\end{column}

\begin{column}{0.5\textwidth}

\end{column}

\end{columns}
    
\end{frame}



\begin{frame}{Ejemplo 3: MOSFET en región lineal}

Proponga una forma de conectar un transistor MOS de manera que la resistencia del canal sea exactamente de $1000\ \Omega$. Asuma $\mu_n C_{ox}' = 100\ \mu{}A/V^2$ y $V_{TH} = 0.4\ V$.
    
\end{frame}


\section{Modelos de simulación}
\begin{frame}{Modelos de Nivel 1, 2 y 3}

Modelo de nivel 1:

\begin{itemize}
	\item Corriente de subumbral = 0
	\item Reemplaza L por L\textsubscript{eff}
\end{itemize}


Modelo de nivel 2:

\begin{itemize}
    \item Saturación de la velocidad de arrastre
    \item Ecuación de corriente de subumbral
    \item Efectos de canal corto y canal angosto, en parte con parámetros de ajuste
\end{itemize}


Modelo de nivel 3:

\begin{itemize}
	\item Modelado de efectos de canal corto, válidos para L hasta de 2 \textmu{}m
	\item Modelado empírico: menos ecuaciones y reducción de tiempo de simulación
\end{itemize}
\end{frame}


\begin{frame}{Modelo BSIM}

\begin{itemize}
	\item BSIM: Berkeley Short-Channel IGFET Model
	\item Modelo más utilizado en simulación
	\item Alrededor de 150 parámetros. Incluye
	\begin{itemize}
		\item Efectos de canal corto
		\item Corriente de compuerta
		\item Efectos de temperatura
		\item Modulación de canal corto
		\item Saturación de velocidad
		\item Corriente de subumbral
		\item Varios modelos de capacitancia
		\item DIBL (drain induced barrier lowering)
		\item Ruido, entre otros
	\end{itemize}
\end{itemize}

\url{http://bsim.berkeley.edu/}

\end{frame}


\begin{frame}{Lecturas recomendadas}

\begin{itemize}
    \item Razavi, B. (2013). Fundamentals of microelectronics. Chapter 6: Physics of MOS transistors, 2nd ed., pp. 270-288, Wiley.
\end{itemize}

\end{frame}

\end{document}
